After the development of the real-space Chern markers detailed in the previous chapter, it was suddenly possible to precisely characterise the topological properties of a band in a wide variety of disordered -- and thus previously somewhat inaccessible -- contexts. This led to a number of works analysing bulk invariants and protected edge modes in amorphous insulators~\cite{mitchell_amorphous_2018,agarwala_higher-order_2020,marsal_topological_2020,costa_toward_2019,agarwala_higher-order_2020,spring_amorphous_2021,corbae_evidence_2019}. However, investigation into amorphous topological physics has been predominantly focused on non-interacting electrons with a few exceptions, for example, to account for the observation of superconductivity in amorphous materials~\cite{buckel_einflus_1954,mcmillan_electron_1981,meisel_eliashberg_1981,bergmann_amorphous_1976,manna_noncrystalline_2022} or very recently to understand the effect of strong electron repulsion in TIs~\cite{kim_fractionalization_2023}. Thus, it is worth asking if there are more exotic physical systems where one might be able to similarly characterise the topological properties of a quantum system without crystalline structure. 
One particularly rich context for finding topological physics is in quantum magnetism, where the properties of amorphous systems have been investigated since the 1960s \cite{aharony_critical_1975,petrakovskii_amorphous_1981,fahnle_monte_1984, kaneyoshi_introduction_1992,plascak_ising_2000,kaneyoshi_amorphous_2018}. Generally research has focussed on classical Heisenberg and Ising models, which are able to realise ferromagnetic and disordered antiferromagnetic, as well as spin glass behaviour~\cite{edwards_theory_1975,coey_amorphous_1978,binder_spin_1986}. 

However, little work has been done to understand whether it is possible to realise a quantum spin liquid (QSL) phase in an amorphous material. There is good intuitive reason why one might avoid this topic. Amorphous magnets often have glassy physics and the problem of finding even the classical ground state is famously difficult. Such systems have an exponentially large number of local minima, and finding the correct one is usually an NP-hard problem \cite{barahona_computational_1982} with new numerical techniques still being developed today \cite{roma_ground_2009,wang_comparing_2015,fan_searching_2023}. Thus, one might reasonably expect that the inclusion of quantum fluctuations would do little to ease the intractability of this problem. If one is searching for QSL phases, looking at amorphous systems would seem to be an even more ill-considered idea. Structural disorder is generally expected to undermine the delicate classical degeneracy necessary for spin liquids to form, instead selecting some arbitrary complex classical state as the ground state and destroying the prospect of forming a long-range entangled ground state.

In the case of a geometrically frustrated candidate QSL, such as the kagome antiferromagnet, the classical frustration is a consequence of lattice structure. Thus, there is little reason to expect that the same model constructed on an arbitrary lattice would inherit any interesting properties of the original frustrated-lattice model. In order to find an amorphous QSL, we turn our attention to exchange frustration. Given an arbitrary lattice, we can ask whether it is possible to construct a set of frustrated couplings such that the system has a QSL ground state. One of the paradigmatic examples of an exchange-frustrated QSL is the Kitaev honeycomb model \cite{kitaev_anyons_2006}. Here, adjacent spins are coupled with three bond-anisotropic Ising interactions in the $x$, $y$ and $z$ directions, which each try to align the spin on the three incommensurate axes. Thus, it is clear that the system has extensive frustration even before considering the geometry of the lattice, making it a promising candidate for realising an amorphous QSL. As we shall review in the following section, Kitaev showed that the model could be exactly solved by transforming to a system of Majorana fermions. In the following years, several instances of Kitaev candidate materials were synthesized~\cite{hermanns_physics_2018,winter_models_2017,trebst_kitaev_2022,takagi_concept_2019} following the suggestion that heavy-ion Mott insulators formed by edge-sharing octahedra may realise dominant Kitaev interactions~\cite{jackeli_mott_2009}. Recently, it has been shown that the Kitaev material Li\textsubscript{2}IrO\textsubscript{3} can be created with an amorphous structure~\cite{lee_lithiair_2019}. In fact, with sufficiently fast cooling, any crystalline material can be made amorphous~\cite{zallen_physics_2008,turnbull_under_1969}, opening the possibility for exploring a wide variety of non-crystalline Kitaev materials, although it must be noted that in such materials the strength of the Kitaev-type interaction and suppression of competing couplings generally depends on the geometric details of the lattice around each ion. Thus, distortions in lattice structure will allow perturbing interactions to be introduced into the Hamiltonian that threaten the stability of the Kitaev QSL phase \cite{rau_trigonal_2014,sela_order-by-disorder_2014,trebst_kitaev_2022}.

By now it is well known that the Kitaev model on any three-coordinated ($z=3$) lattice retains many of the properties that allowed for an exact solution to be constructed in the Honeycomb case. The system retains the conserved plaquette operators and local symmetries that facilitate a mapping onto an effective free Majorana fermion problem in a background of static $\mathbb Z_2$ fluxes~\cite{baskaran_exact_2007,baskaran_spin-s_2008,nussinov_bond_2009,obrien_classification_2016,yao_exact_2007,hermanns_weyl_2015}. However, this neither means that any $z=3$ lattice Kitaev model can be straightforwardly constructed, nor that the QSL properties are obvious. Several obstacles remain. First, the labelling of bonds necessary to create a soluble Hamiltonian can be an NP-complete problem. Second, once the Majorana system has been constructed, determining the ground state out of the exponentially large number of $\mathbb Z_2$ flux sectors is generically hard, since Lieb's theorem~\cite{lieb_flux_1994} -- which defines the ground state flux configuration for the translationally invariant honeycomb -- is not applicable for most lattices. Previous studies have relied on translation and reflection symmetries to reduce the number of sectors that must be checked~\cite{yao_exact_2007,eschmann_thermodynamics_2019,peri_non-abelian_2020}, which cannot be done in an amorphous system. Third, once the ground state flux sector is found, it needs to be determined whether lattice disorder induces a gapless phase~\cite{self_thermally_2019,laumann_disorder-induced_2012,lahtinen_topological_2012,self_thermally_2019} or whether the fermionic spectrum is gapped, possibly with non-trivial topology~\cite{yao_exact_2007}.

In this chapter we introduce the amorphous Kitaev model and establish it as an example of a topologically ordered amorphous QSL. We start in \textsection\ref{sec:honeycomb_model} with a review of the Kitaev model on the Honeycomb lattice. We present in detail the transformation from spins to Majoranas, the role of the $\mathbb Z_2$ gauge field, the role of time reversal symmetry, the ground-state phase diagram and finally the properties of the Majorana excitations. In \textsection\ref{sec:amorphous_kitaev_model} we turn our attention to the amorphous case. Concentrating on random lattices generated via Voronoi tessellation~\cite{mitchell_amorphous_2018,marsal_topological_2020} with $z=3$, we show how to colour the bonds consistently. We find that the presence of plaquettes with an odd-number of sites lead to a chiral QSL with spontaneously broken time-reversal symmetry (TRS)~\cite{yao_exact_2007,chua_exact_2011,chua_exactly_2011,fiete_topological_2012,natori_chiral_2016,wu_-matrix_2009, wang_majorana_2021} and establish via extensive numerics that the ground state $\mathbb Z_2$ flux sector follows a remarkably simple counting rule consistent with Lieb's theorem. We map out the phase diagram of the model and show that the chiral phase around the symmetric point is gapped and characterized by a quantized local Chern number with protected chiral Majorana edge modes. Finally, we discuss the effect of additional bond disorder, showing that the QSL phase is remarkably resistant to local bond-length dependent disorder.